% ─────────────────────────────────────────────────────────────────────────────
Here we describe the background mathematics for the key techniques used in
the BookForMe system: the Elo rating system, Softmax classification, focal loss,
and the expected-score calculation underlying intent confidence thresholding.

% ─────────────────────────────────────────────────────────────────────────────
\section{Elo Rating System}
\label{app:math-elo}

The Elo rating system~\cite{elo1978} assigns each player a numerical rating
$R \in \mathbb{R}$ that reflects their skill level relative to other players.
After each match between a player with rating $R_A$ and an opponent with rating
$R_B$, both ratings are updated as follows.

\subsection{Expected Score}

The expected score $E_A$ for player $A$ against player $B$ is:

\begin{equation}
  E_A = \frac{1}{1 + 10^{(R_B - R_A)/400}}
  \label{eq:elo-expected-full}
\end{equation}

This is a logistic function scaled so that a rating difference of 400 points
corresponds to an expected score of approximately 0.91 (the stronger player wins
roughly 91\% of the time).

\subsection{Rating Update}

After the match, each player's rating is updated by:

\begin{equation}
  R'_A = R_A + K \cdot (S_A - E_A)
  \label{eq:elo-update-full}
\end{equation}

\noindent where:
\begin{itemize}
  \item $K = 32$ is the development coefficient (higher values produce larger
        rating swings; BookForMe uses $K = 32$ for all players, consistent with the
        standard chess Elo implementation for active players).
  \item $S_A \in \{0, 0.5, 1\}$ is the actual outcome score: 1 for a win, 0.5 for
        a draw, 0 for a loss.
  \item $E_A$ is the expected score from Equation~(\ref{eq:elo-expected-full}).
\end{itemize}

\subsection{Matchmaking Window}

The matchmaking search window $\Delta(t)$ expands with wait time $t$ (in seconds)
to guarantee match completion:

\begin{equation}
  \Delta(t) = 200 + 50 \cdot \left\lfloor \frac{t}{60} \right\rfloor
  \label{eq:elo-window}
\end{equation}

A new lobby is created when no existing lobby satisfies
$|R_u - \bar{R}_{\text{lobby}}| \leq \Delta(t)$, where $\bar{R}_{\text{lobby}}$
is the average Elo of current lobby participants.

% ─────────────────────────────────────────────────────────────────────────────
\section{Softmax Classification}
\label{app:math-softmax}

The intent classification head of the DistilBERT model applies a softmax function
over the five logit outputs $z_1, \ldots, z_5$:

\begin{equation}
  P(\text{intent} = k \mid \mathbf{x}) = \frac{e^{z_k}}{\sum_{j=1}^{5} e^{z_j}}
  \label{eq:softmax}
\end{equation}

\noindent The predicted intent class is $\hat{k} = \arg\max_k P(k \mid \mathbf{x})$.

\subsection{Confidence Thresholding}

The dialogue state validation node uses a confidence threshold $\tau = 0.65$:
if $\max_k P(k \mid \mathbf{x}) < \tau$, the agent routes to the Gemini LLM fallback
for intent re-classification rather than proceeding with the low-confidence
DistilBERT prediction.  This prevents the agent from executing incorrect booking
actions based on ambiguous inputs.

% ─────────────────────────────────────────────────────────────────────────────
\section{Focal Loss}
\label{app:math-focal}

Standard cross-entropy loss treats all examples equally:

\begin{equation}
  \mathcal{L}_{\text{CE}} = -\sum_{k=1}^{C} y_k \log \hat{p}_k
  \label{eq:ce-loss}
\end{equation}

For our imbalanced dataset (42\% \texttt{inquiry}, 9\% \texttt{unknown}), standard
cross-entropy causes the model to over-optimise for the majority class.  Focal
loss~\cite{an2022} addresses this by down-weighting easy (well-classified) examples
via a modulating factor $(1 - \hat{p}_t)^\gamma$:

\begin{equation}
  \mathcal{L}_{\text{FL}} = -\sum_{k=1}^{C} y_k \cdot (1 - \hat{p}_k)^\gamma \cdot \log \hat{p}_k
  \label{eq:focal-loss}
\end{equation}

\noindent where $\gamma \geq 0$ is the focusing parameter (we use $\gamma = 2.0$).
When the model is uncertain about a class ($\hat{p}_k \approx 0.5$), the factor
$(1 - 0.5)^2 = 0.25$ reduces the loss contribution by 75\% relative to perfectly
confident examples, focusing gradient updates on the hard minority-class samples.

% ─────────────────────────────────────────────────────────────────────────────
\section{Optimistic Concurrency Control (OCC) Correctness}
\label{app:math-occ}

Let $S_i$ be the state of slot $i$ at time $t$, with $S_i \in \{\text{AVAILABLE},
\text{HOLD}, \text{CONFIRMED}\}$.  Let two concurrent booking requests $r_1$ and
$r_2$ arrive at times $t_1 < t_2$ for the same slot $i$.

In the OCC protocol, each request performs:
\begin{enumerate}
  \item \textbf{Read:} $r_j$ reads $S_i^{(t_j)}$.
  \item \textbf{Verify:} $r_j$ checks $S_i^{(t_j)} = \text{AVAILABLE}$.
  \item \textbf{Write (atomic transaction):} $r_j$ atomically writes
        $S_i \leftarrow \text{HOLD}$ if and only if the document version at write
        time equals the version read in step 1.
\end{enumerate}

By Firestore's serializable transaction isolation guarantee, only one of $r_1, r_2$
can win the atomic compare-and-swap.  Suppose $r_1$ wins.  Then when $r_2$ executes
step 3, the document version has changed (due to $r_1$'s write), and Firestore
returns a transaction conflict error.  $r_2$ then invokes the alternative-slot
query rather than retrying, ensuring exactly-once booking semantics.

%%% Local Variables:
%%% mode: latex
%%% TeX-master: "../report"
%%% End:
